% !TeX root = ../main.tex
\newcommand{\LogitWoeIterations}{30}
\newcommand{\LogitRawIterations}{17}
\newcommand{\LogitMaxVIF}{1.25}
\newcommand{\LogitCondition}{1.83}
\newcommand{\LogitTrainingRate}{6.6844\%}

\begingroup
\setlength{\intextsep}{7pt}
\setlength{\textfloatsep}{10pt}

\subsection{模型估计与诊断}
将第~\ref{sec:feature-selection}~节的 8 项冻结 WOE 输入模型，以坏样本为正类：
\begin{equation}
  Y_i\mid\boldsymbol{x}_i\sim\mathrm{Bernoulli}(p_i),\qquad
  \eta_i=\log\frac{p_i}{1-p_i}=\beta_0+\sum_{j=1}^{8}\beta_jw_j(x_{ij}).
  \label{eq:woe-logistic}
\end{equation}
概率为 $p_i=(1+e^{-\eta_i})^{-1}$。按记录采用条件独立的工作假设，重复与同源限制沿用第~\ref{sec:sample-split}~节；
条件对数几率对 WOE 加性，无需八个特征相互独立。固定 $C=1$，以无类别权重的 L2 惩罚估计
\begin{equation}
  \widehat{\boldsymbol\beta}
  =\arg\min_{\boldsymbol\beta}\left\{
    \sum_{i\in\mathcal T}\bigl[\log(1+e^{\eta_i})-Y_i\eta_i\bigr]
    +\frac{1}{2C}\sum_{j=1}^{8}\beta_j^2\right\},
  \label{eq:penalized-logistic}
\end{equation}
其中 $\mathcal T$ 为训练集，截距不惩罚。采用 \texttt{lbfgs}数值优化算法、容差 $10^{-8}$、最多 2,000 次迭代，
在 \LogitWoeIterations{} 次后收敛。含截距的标准化诊断矩阵秩为 9、条件数为 \LogitCondition，最大 VIF 为 \LogitMaxVIF。
后可根据此惩罚估计得到八个候选变量的WOE Logistic的估计系数。

\begin{table}[!htbp]
  \centering
  \caption{八变量 WOE Logistic 的估计系数}
  \label{tab:logistic-coefficients}
  \footnotesize
  \setlength{\tabcolsep}{8pt}
  \renewcommand{\arraystretch}{1.02}
  \begin{tabular}{lrr}
\toprule
项 & 系数 $\widehat\beta$ & $\exp(\widehat\beta)$ \\
\midrule
截距 & -2.60255 & 0.0741 \\
额度使用率 & -0.61583 & 0.5402 \\
年龄 & -0.40832 & 0.6648 \\
30–59 天逾期次数 & -0.51178 & 0.5994 \\
负债收入比 & -0.76611 & 0.4648 \\
90 天及以上逾期次数 & -0.53753 & 0.5842 \\
房地产贷款数 & -0.67415 & 0.5096 \\
60–89 天逾期次数 & -0.35875 & 0.6986 \\
家属人数 & -0.28302 & 0.7535 \\
\bottomrule
\end{tabular}

  \begin{minipage}{\linewidth}
    \footnotesize
    注：非截距行的指数表示其他 WOE 固定时，该变量增加 1 个 WOE 单位的坏样本几率比；
    截距项的指数表示所有 WOE 为 0 时的坏/好几率。
  \end{minipage}
\end{table}

表~\ref{tab:logistic-coefficients} 八项系数均为负：WOE 增大时条件坏样本几率下降。
以负债收入比为例，当其他七项 WOE 不变时，负债收入比的 WOE 增加 1 个单位，
坏样本几率变为原来的 0.5402 倍，即降低约 45.98\%。
年龄系数的单位也是 WOE，并非每增加一岁的效应，截距不代表平均客户风险。

\subsection{具体档位的条件几率比}
同一变量由档位 $a$ 变为 $b$，其余七项输入固定时，
\begin{equation}
  \mathrm{OR}_{b:a}
  =\frac{p_b/(1-p_b)}{p_a/(1-p_a)}
  =\exp\!\left\{\widehat\beta_j\bigl(w_{j,b}-w_{j,a}\bigr)\right\}.
  \label{eq:conditional-bin-odds-ratio}
\end{equation}

以第~\ref{sec:scorecard}~节的低风险人工输入为参照。负债收入比分别设为 0.25 与 2，严重逾期次数分别设为 0 与 1。
分别只改设负债比或严重逾期档位，结果见表~\ref{tab:logistic-odds-examples}。
DebtRatio 的 WOE 差为 $-0.65815$，与系数 $-0.76611$ 结合得到 $\mathrm{OR}\approx1.656$；
严重逾期的 WOE 差为 $-2.35001$，与系数 $-0.53753$ 结合得到 $\mathrm{OR}\approx3.537$。
几率倍数不等于概率倍数；概率变化依赖参照特征，表示模型情景比较，不解释为因果效应。

\begin{table}[!htbp]
  \centering
  \caption{固定其余特征后的档位比较}
  \label{tab:logistic-odds-examples}
  \footnotesize
  \setlength{\tabcolsep}{6pt}
  \begin{tabular}{llrrr}
\toprule
变量 & 档位变化 $a\to b$ & 条件几率比 & $p_a$ & $p_b$ \\
\midrule
负债收入比 & $(0,0.4]\to(1,10]$ & 1.656 & 1.13\% & 1.85\% \\
90 天及以上逾期次数 & $0\to1$ & 3.537 & 1.13\% & 3.88\% \\
\bottomrule
\end{tabular}

\end{table}

\subsection{两个基准与同样本比较}
\paragraph{常数概率基准}
对所有客户预测训练集坏样本率 $\widehat p_0=6{,}016/90{,}000=\LogitTrainingRate$，
不从测试标签重估先验。
\paragraph{原始字段 Logistic 基准}
同样八个字段中，年龄保留原值，其余取 $\log(1+x)$。
96/98 对应普通计数置 0，以单个“任一逾期特殊”指示保留状态，共 9 个设计量，避免三个相同标记重复输入。
按训练均值和总体标准差缩放后冻结，使用同一 $C$、求解器及无权重设置，
17 次后收敛，矩阵满列秩且未检出分离。

两模型参数和规则先冻结，随后才打开同一批 30,000 条测试记录（坏样本 2,005 条），结果见表~\ref{tab:logistic-comparison}。
两个 Logistic 模型采用相同的八个原始字段，分别使用数值变换后的原值表示和分箱后的 WOE 表示。
测试结果显示，WOE 流程的 AUC 提高约 0.0235，KS 和 AP 同时提高，
对数损失与 Brier 分别降低约 4.53\%和 2.91\%，
表明该流程在本次实验中改善了风险排序与整体概率预测质量。
这与分箱能够表达不同取值区间的风险差异相一致，例如 DebtRatio 的非单调关系难以通过单一对数变换充分表达。

\paragraph{不平衡与准确率}
训练沿用约 $13.96:1$ 的好坏比，各记录等权。给定模型预测的坏样本概率 $\widehat p_i$，
记分类阈值为 $t$，由此定义预测标签
\[
  \widehat Y_i(t)=
  \begin{cases}
    1,\ \text{标记为高风险客户}, & \widehat p_i\geq t,\\
    0,\ \text{不标记为高风险客户}, & \widehat p_i<t.
  \end{cases}
\]

若误报与漏报代价相同且预测概率可信，标记和不标记的期望成本分别与
$1-\widehat p_i$ 和 $\widehat p_i$ 成正比。当 $1-\widehat p_i\leq\widehat p_i$ 时选择标记，
得到等代价对照阈值 $t_0=0.5$。（即此时考虑误报好客户成本与漏报坏客户成本相同）

以坏样本为正类，将 WOE 模型的 $\widehat Y_i(t_0)$ 与测试集实际标签逐一比较，
得到正确检出 $TP=364$、漏报 $FN=1{,}641$、误报 $FP=244$、正确识别好客户 $TN=27{,}751$。
准确率与坏样本召回率的定义及代入结果分别为
\[
  \begin{gathered}
    \mathrm{Accuracy}(t)=\frac{TP(t)+TN(t)}{N},\qquad
    \mathrm{Recall}_{\mathrm{bad}}(t)=\frac{TP(t)}{TP(t)+FN(t)},\\[0.3ex]
    \mathrm{Accuracy}(t_0)=\frac{364+27{,}751}{30{,}000}=93.72\%,\qquad
    \mathrm{Recall}_{\mathrm{bad}}(t_0)=\frac{364}{2{,}005}=18.15\%.
  \end{gathered}
\]

在该阈值下，好样本识别率为 99.13\%。常数基准的预测概率低于 $t_0$，因而将所有人判为好样本，
准确率等于好样本占比 93.32\%，坏样本召回率为 0。好样本数量占优，使总体准确率主要反映多数类的识别表现。

\endgroup
